\chapter{Psychological theories of syllogistic reasoning}
\label{sec:p19-rationality-3-3}

Syllogistic reasoning supplies a small domain in which a cognitive theory
must explain both conspicuous errors and successful deductions. Khemlani's
2021 treatment does not identify a person's most frequent answer with a valid
conclusion. It asks how people interpret quantified premises, construct a
representation, choose a conclusion and sometimes correct it. Its central
proposal is to integrate biased, probabilistic construction with deliberation
over discrete representations rather than choose between heuristics and
deliberation as exclusive explanations
\cite[pp.~225--236]{Khemlani2021Syllogistic}.

The three objects of comparison are a specified logical consequence, a
theory's predicted response distribution and the responses actually recorded.
The first requires a semantic argument, the second an algorithm and parameter
specification, and the third a task and evidence record. Their interfaces
follow \cref{def:p19-benchmark,def:p19-evidence-record}.
\Cref{sec:p02-syllogistic-inference} supplies a modern formal calculus,
not a psychological explanation of selecting one answer among many.

\section{The task, its response space and its existence convention}
\label{sec:p19-r33-task}

\begin{hbdefinition}{def:p19-r33-categorical-task}{The chapter's categorical task notation}
In this chapter's task scope, lower-case \(a,b,c\) name unary properties;
\(b\) is the middle property and \(a,c\) are the end properties.
For a nonempty universe \(U\), let \(a^U,b^U,c^U\) be their extensions.
The task abbreviations have the following declared meanings:
\[
\begin{array}{c|l|l}
Aab&\text{all }a\text{ are }b&a^U\subseteq b^U\\
Iab&\text{some }a\text{ are }b&a^U\cap b^U\ne\varnothing\\
Eab&\text{no }a\text{ is }b&a^U\cap b^U=\varnothing\\
Oab&\text{some }a\text{ are not }b&a^U\setminus b^U\ne\varnothing.
\end{array}
\]
The capital letters are sentence-mood codes, not extensions of properties.
The ordered premise figures used by the source are
\[
\begin{array}{c|cccc}
\text{figure}&1&2&3&4\\
\text{first premise terms}&a\!:\!b&b\!:\!a&a\!:\!b&b\!:\!a\\
\text{second premise terms}&b\!:\!c&c\!:\!b&c\!:\!b&b\!:\!c.
\end{array}
\]
This is premise-only figure numbering. It is not silently interchangeable
with a scheme that fixes the order of the conclusion.
\end{hbdefinition}

There are \(4\times4\times4=64\) ordered pairs of premise forms.
The usual requested answers are the eight categorical sentences
\(Aac,Iac,Eac,Oac,Aca,Ica,Eca,Oca\), or \emph{no valid conclusion}
(NVC). Thus the 64 premise tasks, the 512 possible premise-pair/categorical-
answer combinations and the nine response categories are different counts.
In a generation task, participants formulate answers; in an evaluation or
selection task, experimenters supply candidates. The latter can prompt
answers a person would not otherwise produce
\cite[pp.~226--227]{Khemlani2021Syllogistic}.

\begin{hbconvention}{def:p19-r33-existence}{Keep existence information in the benchmark}
The source's 27-pairs-with-a-definite-conclusion and 37-NVC-pairs division
uses the psychological syllogistic benchmark in which all three relevant
classes have members. The underlying meta-analysis states that assumption
explicitly and records instructions or linguistic conventions supplying
existence \cite[pp.~3--4 of the advance-online text]
{KhemlaniJohnsonLaird2012SyllogismMeta}.
Where that benchmark is used below, add
\(a^U,b^U,c^U\ne\varnothing\). Elsewhere the displayed set semantics
allows empty extensions, as does P2. An existential conclusion requires an
actual witness or a separately stated nonemptiness premise.
NVC denies a licensed answer within the eight specified end-term forms;
it does not deny all logical consequences of the premises.
\end{hbconvention}

The chapter's first illustration is: all women are designers; some employees
are not women; what follows? With \(a=\)designer, \(b=\)woman,
\(c=\)employee, it is \(Aba,Ocb\). The source contrasts a frequent
erroneous existential-negative answer with the less frequent NVC answer.
The example is not a deduction from occupations or sex categories as facts;
the stated premises are hypothetical inputs to the task.

\begin{hbexample}{def:p19-r33-easy-hard}{The easy and difficult designer examples}
For \(a=\)designer, \(b=\)animator and \(c=\)undergraduate,
the source's easy example is \(Iab,Abc\); its conclusion \(Iac\)
uses the witness in the first premise. The difficult example is
\(Eab,Abc\). Its intended conclusion \(Oca\) uses an animator
witness, not merely the absence of designer-animator overlap.
A designer can also be an undergraduate without being an animator, so
the tempting \(Eac\) conclusion does not follow. The source reports a
large difference in response patterns between these tasks
\cite[pp.~226--228]{Khemlani2021Syllogistic}.
\end{hbexample}

\begin{hblemma}{thm:p19-r33-witness}{proved=annex:p19.r33.witness}{The exact witness contrast}
Under \cref{def:p19-r33-categorical-task}, \(Iab,Abc\) entails
\(Iac\). The premises \(Eab,Abc\), together with
\(b^U\ne\varnothing\), entail \(Oca\), but do not entail
\(Eac\). Without the nonemptiness condition they need not entail
\(Oca\), even if the universe and both end classes are nonempty.
\end{hblemma}

\begin{hblemma}{thm:p19-r33-nvc}{proved=annex:p19.r33.nvc}{No definite conclusion in two source examples}
Both \(Aba,Ocb\) and \(Aab,Obc\) admit nonempty extensions
of all three properties while entailing none of the eight end-term
categorical sentences. Thus NVC is warranted for both tasks under the
declared response benchmark. The proof provides complete countermodels
covering every candidate, not merely a failure to find one conclusion.
\end{hblemma}

\section{Phenomena that an account must explain}
\label{sec:p19-r33-phenomena}

The chapter opens with Störring's 1908 study: four volunteers were observed
solving a battery in a dark room, with verbal answers, times and bodily
behavior recorded. Some answers appeared immediate; others elicited reports
of deliberation, imagery or diagrams. Khemlani treats this as a historical
origin, while noting the arbitrary design and absence of a modern
quantitative analysis. These reports do not independently identify an
underlying process or support general population frequencies. The chapter's
later neuroimaging and clinical-testing examples explain why syllogisms
became a microcosm; they are not a validated diagnosis of an individual's
rationality \cite[pp.~225--226,233--234]{Khemlani2021Syllogistic}.

Six kinds of variation motivate a process account. Problems differ in
difficulty; individuals differ in accuracy; term order changes the
direction and frequency of conclusions; believable contents affect
acceptance; understanding quantifiers develops; and people vary or change
their strategies. The figural effect is especially informative: two
logically equivalent premises can produce different answer directions
without changing which conclusions are valid. A response bias cannot
therefore be explained by a list of valid syllogisms alone. Conversely,
belief bias requires a task with meaningful contents and cannot be inferred
from an abstract-item fit \cite[Table~3.3.1, p.~228]{Khemlani2021Syllogistic}.

The source groups earlier theories by their explanatory emphasis. Heuristic
accounts use atmosphere, matching, illicit conversion or probabilistic
heuristics. Simulation accounts represent a possible situation and may
inspect alternatives. Mental-rule accounts construct proofs under
psychological control constraints. These are not interchangeable: an
atmosphere or matching account mainly constrains the mood of an answer,
whereas a representation with an intrinsic order can also constrain its
direction. A proof-based account needs an account of how a candidate is
generated and how NVC is recognized. The chapter sets aside detailed
rule-theory exposition after stating this difficulty; it does not show
that every possible cognitive rule-and-control account is impossible
\cite[pp.~228--229]{Khemlani2021Syllogistic}.

\begin{hbexample}{def:p19-r33-matching}{Matching two existential premises}
The source's example is: some designers are women; some women are employees.
A matching strategy proposes some designers are employees, retaining the
existential mood. But the two premise witnesses need not be the same
woman. An initial representation that merges them supports that answer;
separating them can refute it. The meta-analysis records 61\% \(Iac\)
responses and 33\% NVC responses for \(Iab,Ibc\), rounded aggregate
response proportions rather than population probabilities
\cite[Table~6, p.~22 of the advance-online text]
{KhemlaniJohnsonLaird2012SyllogismMeta}.
\end{hbexample}

\begin{hblemma}{thm:p19-r33-split}{proved=annex:p19.r33.split}{Distinct witnesses block the matching inference}
\(Iab,Ibc\) is consistent with \(Eac\), and also consistent
with \(Iac\), even with all three classes nonempty.
Thus its premises do not entail either of those incompatible answers.
Splitting one triple-property individual into two individuals can preserve
both existential premises while refuting \(Iac\).
\end{hblemma}

These contrasts motivate the chapter's objection to a false dichotomy.
A heuristic-only account has difficulty explaining a person's corrected
deduction; an unrestricted deliberative account can predict too many
responses and too much search. A viable integration must restrict how an
initial representation is built, what it makes accessible, when it is
revised and how a revised representation changes the reported answer.
Naming fast and slow systems without specifying those operations is not
yet an explanation. Sequential proposal-and-check processing is one
architecture; the broader dual-process view also permits parallel,
interacting operations \cite[pp.~229--230]{Khemlani2021Syllogistic}.

\section{What the 2012 theory comparison actually measured}
\label{sec:p19-r33-meta}

The underlying meta-analysis surveyed twelve theories but made a quantitative
comparison of seven with sufficiently explicit predictions: atmosphere,
matching, conversion, PHM, PSYCOP with its implicature assumptions, verbal
models and the then-current mental-model implementation. Omitted accounts
lacked the required all-64 predictions or were not proposed as complete
error-generating accounts. Their omission is not a measured loss to the
seven compared models. The authors searched through 2011 and selected
studies requiring self-generated conclusions to all 64 pairs, excluding
studies confined to conclusions in one Scholastic order
\cite[pp.~19--21 of the advance-online text]
{KhemlaniJohnsonLaird2012SyllogismMeta}.

\begin{longtable}{@{}p{0.38\linewidth}p{0.08\linewidth}p{0.41\linewidth}@{}}
\caption{The six datasets in the 2012 comparison}\label{tab:p19-r33-meta-samples}\\
\hline
Dataset & \(N\) & Reported population and instruction \\
\hline
\endfirsthead
\caption[]{The six datasets in the 2012 comparison (continued)}\\
\hline Dataset & \(N\) & Reported population and instruction \\
\hline
\endhead
\hline\multicolumn{3}{r}{\emph{continued on next page}}\\
\endfoot
\endlastfoot
Johnson-Laird--Steedman 1978, first session & 20 & Teachers College,
Columbia students; deduce with absolute certainty. \\
Same experiment, second session & 20 & The same 20 students retested;
not a new independent cohort. \\
Johnson-Laird--Bara 1984, Experiment~3 & 20 & Milan undergraduates;
what follows necessarily. \\
Bara et al. 1995, adolescent group & 20 & Florence high-school students
aged 14--15; what follows necessarily. \\
Bara et al. 1995, adult group & 20 & Florence undergraduates;
what follows necessarily. \\
Roberts et al. 2001 & 56 & Florida undergraduates;
what follows necessarily. \\
\hline
\end{longtable}

These are the meta-analysis's provenance records, not a claim that all
underlying participant-level protocols have been reanalyzed here. The
six nominal sample counts sum to 156 session records, not 156 distinct
people: the Columbia sessions repeat participants. Common laboratories,
convenience recruitment and differences in instructions limit independence
and transfer. The source weights comparisons by study sample size and
reports confidence intervals and heterogeneity checks; this does not make
its six datasets six randomized, mutually independent trials
\cite[Table~5, p.~21 of the advance-online text]
{KhemlaniJohnsonLaird2012SyllogismMeta}.

The metric is deliberately not participants' logical accuracy. A theory
predicts which answers occur often enough, and should also exclude answers
that do not. The analysis declares a \(1/9\) response-chance reference
and describes a 16\% frequency cutoff. Its prose says ``at least 16\%,''
while Table~8's footnote uses strict greater/less-than 16\%; an exact
reproduction must resolve that threshold-boundary convention from the
coding, not from rounded tables. The combination of hits and correct
rejections penalizes a theory that predicts every answer. It compares
possible-response sets, which can disadvantage an account designed to
predict detailed probabilities instead
\cite[pp.~19--21,25 of the advance-online text]
{KhemlaniJohnsonLaird2012SyllogismMeta}.

The then-current mental-model account had hits 95\% (reported 95\% CI
91--99) but correct rejections 74\% (70--78): its prediction set was too
liberal. Verbal models had 84\% combined correct predictions (80--89),
conversion 83\% (80--86), and mental models 78\% (75--81). The verbal-
model combined result had marginal heterogeneity, \(I^2=52\%\),
\(Q=10.69,p=.06\); the mental-model hits analysis had the reported
publication-bias regression \(p=.04\). These estimates belong to the
source's thresholded comparison of particular implementations, not a
ranking of rational capacities or a theorem about all future model theories
\cite[Table~8, p.~25 of the advance-online text]
{KhemlaniJohnsonLaird2012SyllogismMeta}.

The source's heatmap is a visualization of these response patterns, not a
new dataset. Its white cells include percentages at or below 16\%, so white
does not mean nobody made the response. The exact underlying table gives
88\% valid answers for \(Iab,Abc\) and 3\% \(Oca\) answers for
\(Eab,Abc\); the latter instead drew 71\% \(Eac\) answers.
The chapter's wider 1--90\% difficulty range concerns other items in this
bounded collection. No unreported interval for those item proportions is
invented \cite[pp.~226--228]{Khemlani2021Syllogistic}
\cite[Table~6, p.~22 of the advance-online text]
{KhemlaniJohnsonLaird2012SyllogismMeta}.

\section{Hattori's probability sampling account}
\label{sec:p19-r33-psm}

The chapter calls Hattori's account the \emph{probability sampling model}
(PSM). Hattori's original published paper calls its overall account the
\emph{probabilistic representation model} (PRM). They refer here to the
same 2016 presentation, not two separately validated theories. Its
probability prototype model (PPM) precedes a sample mental model (SMM).
Euler-style regions describe combinations of membership, but Hattori
explicitly permits equivalent probability-equipped representations; the
geometry is not an obligatory internal picture
\cite[pp.~229--231]{Khemlani2021Syllogistic}
\cite[pp.~300--304]{Hattori2016PRM}.

\begin{hbmethod}{def:p19-r33-psm-procedure}{Prototype, finite sample and ordered tests}
\begin{enumerate}
\item Interpret each premise using a standard representative relation and
combine them with minimal further constraints, including the model's
conditional-independence assumption where applicable. Assign probabilities
to the eight membership regions using the occurrence and coverage
parameters. These choices add psychological assumptions beyond the premises.
\item Draw a small independent sample of membership triples from that
distribution. Repeated draws can occupy the same region, so the number
of draws is not the number of distinct property patterns represented.
\item Test the four end-term moods against the sampled individuals. The
least informative premise selects the start of a cyclic test order, based
on \(A,I,E,O\). A max-heuristic parameter controls probabilistic
acceptance of a passing candidate; failure to accept candidates can lead
to NVC. This acceptance stage must not be omitted from the full algorithm.
\end{enumerate}
In this local prototype scope, \(S,M,P\) name the three membership
events, \(x=P(M)\) is middle-event occurrence, \(c\) is coverage,
\(n\) is sample size and \(u_A,u_I,u_E,u_O\) are acceptance
parameters. Here the \(P\) in \(P(M)\) is the probability operator,
not the third event. These roles are local to this model and distinct
from the lower-case task-property notation.
\end{hbmethod}

The chapter also illustrates sampling from a single premise, all designers
are women. In its declared two-region illustration, both designer-women
and women who are not designers have equal probability. Four example
draws alternate those two types, with repeated types retained as separate
tokens. The universal premise rules out designers who are not women;
it does not itself give the two allowed types equal probabilities or
fix their sampled frequencies. Those are additional representational
choices \cite[p.~230]{Khemlani2021Syllogistic}.

The source's diagram uses all \(P\) are \(M\) and some \(S\)
are \(M\). A prototype allows overlap between \(S,P\), but a small
sample may miss it. The diagram then tests \(I\), finds no sampled
overlap and accepts \(E\). That is truth in the selected finite sample,
not truth in every premise-satisfying situation. For \(x=.20,c=.90\),
Hattori's exact prototype table assigns overlap mass \(cx^2=.036\);
the reproduced figure rounds that region to .04. Neither number is a
measured fraction of participants. With seven independent draws, the
overlap-miss probability is \(.964^7\), approximately .774, conditional
on that prototype and its sampling assumptions
\cite[p.~231, Figure~3.3.2]{Khemlani2021Syllogistic}
\cite[pp.~300--304, Table~2]{Hattori2016PRM}.
An existential premise can similarly have positive witness mass in the
prototype while a finite unconditioned sample contains no witness to it.
Thus consistency of the prototype with the premises does not guarantee
that every sample is itself a complete model satisfying their existential
requirements; a conditioning or repair step would need to be specified.

\begin{hblemma}{thm:p19-r33-sampling}{proved=annex:p19.r33.sampling}{What finite sampling and the categorical tests establish}
For a nonnegative integer \(n\) of independent draws from a distribution
with overlap mass \(q\),
the probability of missing overlap in \(n\) draws is \((1-q)^n\).
For \(0<q<1\) and finite \(n\), it is positive: failure to sample a
witness does not prove its impossibility. Under the displayed set semantics,
\(A\) and \(O\), and \(I\) and \(E\), are complementary
sample tests. Therefore a procedure testing all four and accepting every
passing mood must produce a categorical answer, not NVC. This conclusion
does not apply to Hattori's additional probabilistic acceptance gates.
\end{hblemma}

Hattori initially fixes seven draws and explores sizes four through twelve
against eight previous datasets. Six minimizes RMSD, while seven maximizes
correlation. This is model fitting, not a measurement that a participant's
working memory contains exactly six or seven individuals. The chapter
reports comparable performance against PHM and a parameterized mental-model
account, then emphasizes the ability to alter contents through prototype
probabilities \cite[p.~230]{Khemlani2021Syllogistic}
\cite[p.~312, Figure~14]{Hattori2016PRM}.

Its Frenchmen/wine-drinkers/Italians example has the same \(Iab,Ibc\)
form as \cref{def:p19-r33-matching}. A low assumed French--Italian overlap
reduces the predicted affirmative answer. The change can explain an effect
of background knowledge without making \(Iac\) logically valid or
demonstrating that citizenship categories are mutually exclusive.
The source's two underlying experiments more narrowly manipulate prevalence
and inclusion: rare fictitious blood type, pollen allergy and coughing
were used to alter a default set-size balance assumption. They are not
two new experiments on the French--Italian example
\cite[p.~230]{Khemlani2021Syllogistic}
\cite[pp.~313--314]{Hattori2016PRM}.

Experiment~1 recruited 89 Ritsumeikan undergraduates (reported mean age
20.2, SD 1.0), excluded one instruction failure, and used twelve selected
syllogisms in randomized booklet order with reverse-order versions assigned
randomly. Experiment~2 recruited fifty Kanazawa undergraduates (mean 20.8,
SD 1.0) for voluntary classwork and used 32 computerized tasks; records for
\(IA2,II1,II2,II4\) were lost to a programming error. The fit summaries favored
PRM by very small AIC differences from PHM, 8.080 versus 8.170 and 10.374
versus 10.422, respectively. The fitted coverage values .805 and .779 were
compared to earlier datasets, not to a simultaneous randomized neutral-content
control group. Neither an AIC ordering nor an estimated parameter proves
the prototype is a person's actual representation
\cite[pp.~313--315, Tables~5--6]{Hattori2016PRM}.

The chapter's criticisms are specific: the selected two-premise construction
does not establish arbitrary-premise scaling, and a fixed sequence of sample
tests does not reconstruct the strategies people use to search for
counterexamples. Its blanket claim that PSM cannot produce NVC needs a
qualification: Hattori's full algorithm explicitly includes NVC via
probabilistic acceptance/rejection. That is not the model-revision account
of \emph{why} a person recognizes that no conclusion is guaranteed.
\Cref{thm:p19-r33-sampling} explains the narrower limitation of the
un-gated tests. Computational tractability, a response-category output and
a cognitively substantiated deliberative process are three different
achievements \cite[p.~230]{Khemlani2021Syllogistic}
\cite[p.~304, Figure~8]{Hattori2016PRM}.

\section{mReasoner: construction, scanning and counterexamples}
\label{sec:p19-r33-mreasoner}

mReasoner begins with three commitments: models represent discrete
possibilities; their structure is iconic while permitting abstract negation
symbols; and rapid construction/scanning interacts with slower revision.
Unlike the PPM route, it builds a model directly from interpreted premises.
A canonical model of all designers being women represents designers who
are women; fuller possibilities also allow women who are not designers.
Biased construction can omit the latter without denying their semantic
possibility \cite[pp.~231--233]{Khemlani2021Syllogistic}.

\begin{hbdefinition}{def:p19-r33-mreasoner-parameters}{Four parameters in the 2016 implementation}
In this implementation's scope, \(\lambda>0\) controls the size
distribution of an initial model; \(\varepsilon\in[0,1]\) controls
drawing from fuller rather than canonical possibilities;
\(\sigma\in[0,1]\) is the probability of engaging counterexample
search; and \(\omega\in[0,1]\), conditional on successful search,
is the probability of weakening and testing the conclusion rather than
immediately responding NVC. The size variable describes model tokens,
not recruited participants. An inferred parameter value is not an
independent observation of the associated psychological process
\cite[pp.~2--3]{KhemlaniJohnsonLaird2016IndividualSyllogisms}.
\end{hbdefinition}

The source chapter's schematic highlights size, noncanonical construction
and revision; the detailed implementation separates engagement from
reconcluding. The final journal article, first published online in 2021
and assigned to a 2022 issue, calls these parameters size, atypicality,
search and reconclude. It explicitly uses a left-truncated Poisson size
distribution. A truncation or minimum-token rule matters; a bare Poisson
law gives positive mass to zero and does not itself supply a nonempty
model. The journal presentation is a qualified related version, not the
unchanged 2021 chapter or the unchanged 2016 simulation
\cite[pp.~295--297]{KhemlaniJohnsonLaird2021OnlineProperties}.

\begin{hbmethod}{def:p19-r33-mreasoner-procedure}{The specified two-pipeline mechanism}
\begin{enumerate}
\item Parse the premises and construct a small, biased model satisfying
their interpreted constraints. Typical and fuller property patterns are
sampled under the model-construction parameters.
\item Scan the model for an initial conclusion. The order in which end
properties enter the representation constrains conclusion direction;
the moods of premises constrain its polarity and particularity.
This selection mechanism is a source of both successful and invalid answers.
\item If search is engaged, alter the model while preserving the premises:
add individuals, split an individual's properties between individuals,
or move properties between individuals. Test whether an alternative
model falsifies the initial conclusion.
\item On finding a counterexample, either report NVC or attempt a weaker
conclusion and search again. Retain the original model's constraints
as well as the alternative's: fitting a conclusion to just the newest
model is not a validity test.
\end{enumerate}
The journal implementation applies its three revision operations together
because their actual order in participants had not been established
\cite[pp.~297--298]{KhemlaniJohnsonLaird2021OnlineProperties}.
\end{hbmethod}

The chapter's principal illustration is all designers are women; some women
are not employees, or \(Aab,Obc\). A canonical representation may
include designer-women lacking an employee token and suggest \(Oac\).
An alternative puts every designer among employees while locating the
required nonemployee woman outside the designer class. That alternative
preserves both premises and refutes \(Oac\). The nonemployee status
must be assigned explicitly in a complete formal countermodel; a blank
cell in an incomplete token diagram is not automatically a negative literal.
\Cref{thm:p19-r33-nvc} establishes the full NVC benchmark for this
example. Finding one counterexample refutes one candidate, whereas NVC
requires ruling out all eight requested candidates
\cite[p.~232]{Khemlani2021Syllogistic}.

\begin{hbtheorem}{thm:p19-r33-finite-check}{proved=annex:p19.r33.finite-check}{A finite exhaustive benchmark, not a finite cognitive sample}
For the three-property A/I/E/O language and the optional nonemptiness of
its named classes, any countermodel to a proposed categorical conclusion
has a countermodel with at most eight individuals. Consequently,
enumerating nonempty sets of the eight complete membership types can
decide whether any of the eight requested conclusions follows.
This does not show that mReasoner or a participant performs the exhaustive
enumeration, or that a short random sample is complete.
\end{hbtheorem}

The finite benchmark distinguishes complete coverage from a successful
search in one case. A model-revision routine must preserve premise truth;
failure to discover a counterexample only warrants validity if its search
space and procedure are exhaustive. Likewise \(Aac\) can be weakened
to \(Iac\) only with an \(a\)-witness, and \(Eac\) to \(Oac\)
only with one. The task's existence convention supplies these bridges;
the premise-free set semantics does not.

\section{Individual fits, changing strategies and unresolved evidence}
\label{sec:p19-r33-individual}

Khemlani and Johnson-Laird's 2016 conference paper reanalyzed the available
responses of twenty Columbia students from the 1978 experiment. It did not
recruit a new twenty-person sample. An exploratory clustering procedure
identified three groups, then a grid of 1,512 combinations of four
parameters generated 100 simulated repetitions of each of 64 tasks per
setting. Best-fitting subgroup correlations were .74, .82, .90; individual
best-fit correlations averaged .70 and ranged .54--.87. The subgroups
were interpreted as intuitive, intermediate and deliberative under the
parameter definitions. These are fits selected using the same responses,
not held-out classifications of stable psychological types; simulated
repetitions are not extra human participants
\cite[pp.~3--6]{KhemlaniJohnsonLaird2016IndividualSyllogisms}.

The later journal version changes the fitted subgroup settings and reports
different best-fit ranges; its table of 11 simulations also separates
generation, necessity, possibility and consistency tasks. That expanded
coverage does not make one task's result another task's replication.
The 2016 paper admits a wrong prediction for \(Oba,Acb\); the journal
version admits overprediction of NVC and failures for some negative
individual-property inferences. The chapter's favorable fit claim retains
these conflicts, rather than implying that every inference is explained
\cite[p.~6]{KhemlaniJohnsonLaird2016IndividualSyllogisms}
\cite[pp.~298--301]{KhemlaniJohnsonLaird2021OnlineProperties}.

Ragni and colleagues' 2018 reanalysis follows those same twenty students
across two individually administered sessions about a week apart. Each
session used all 64 pairs in a different randomized order, occupational
contents, timed spontaneous conclusions, and instructions to be accurate,
quick and certain. Participants did not know at the first session that
they would be retested. Reported logical accuracy rose from 58\% to 68\%;
19 of 20 improved (reported binomial \(p<.0001\)). Their average number
of changed answers was 27.5 of 64, with reported SD 7.25. There was no
randomized no-retest control, no retained item-order data for the two
sessions, and no assurance against studying reasoning between sessions.
The evidence rejects exact within-person response stability in this sample;
it does not by itself identify practice as the sole cause or show transfer
to everyday reasoning \cite[pp.~2--5]{RagniEtAl2018LogicalWithoutFeedback}.

The 2018 model fits used another 1,512-setting grid, then 1,000 simulated
participants per fitted week. Search and weakening parameters increased
in the fitted description, but those changes are not direct measurements
of counterexample search. The aggregate response correlations .84/.85
and the individual-fit summaries 60\%/65\% have different denominators
and meanings. No confidence interval for the 10-percentage-point human
accuracy change is reported in that selected record, and none is supplied
here. This practice analysis and the 2016 clustering reuse the same human
data; they are not independent confirmations in new populations
\cite[pp.~3--5]{RagniEtAl2018LogicalWithoutFeedback}.

The chapter also cites Ragni--Dames--Brand--Riesterer 2019 on NVC. Its
exact proceedings identity is known, but the full original body needed
to qualify recruitment, exclusions, mixed-model estimates and uncertainty
is not available in the present evidence record. The chapter's qualitative
appeal to that work therefore remains a source-attributed motivation,
not a reconstructed empirical test. Neither its indexed abstract nor a
different NVC paper supplies the missing record
\cite[p.~228,230,236]{Khemlani2021Syllogistic}
\cite{RagniEtAl2019NothingFollows}.

\section{Comparison, rationality and task lookup}
\label{sec:p19-r33-lookup}

Both integrated accounts sample biased representations and implement
explicit procedures. They are not an unrestricted union of every older
theory's answers: that would hide incompatible commitments and predict
too much. PSM's prototype introduces calibrated membership probabilities;
mReasoner directly builds discrete models and can revise them in search
of counterexamples. The chapter praises their shared ability to bridge
heuristic and deliberative processing, while identifying complementary
limitations \cite[pp.~233--234]{Khemlani2021Syllogistic}.

mReasoner's ability to accept arbitrary numbers of premises is an
architectural claim, not an empirical proof that every longer task is
explained. The chapter explicitly says that its syllogistic belief-bias
simulation and its account of learning/reasoning interaction were not yet
supplied. A sentential semantic-modulation prototype is not a completed
syllogistic contents explanation. The broader closing list of desirable
coverage must be read with these exclusions, not as their cancellation
\cite[pp.~232--234]{Khemlani2021Syllogistic}.

The P2 calculus proves consequence under declared interpretation and
existence assumptions. The P19 mental-model chapter explains a general
representational proposal; this unit adds quantified-task construction,
response selection and specific simulation comparisons. O'Brien's natural
logic in \cref{sec:p19-rationality-3-2} is an actual rival cognitive
account of conceptual-role interpretation, control and pragmatic
integration, not merely formal rules without a parser.
The PHM/PRM comparison in \cref{sec:p19-rationality-4-5} additionally
retains the generalized-quantifier experiments in which PRM was inferior
to PHM. The present traditional-quantifier content fits cannot overwrite
that result or establish a universal model winner.

The chapter's final rationality argument is a proposal about integration.
People can draw guaranteed conclusions, likely conclusions and mistaken
conclusions; probabilistically constrained construction can lead to
discrete models that support either intuition or deduction. That makes
theories combining those operations promising candidates for overcoming
a theoretical impasse. It does not turn every sampled conclusion into a
normatively justified response, prove that deduction is dispensable,
or validate syllogistic accuracy as a general clinical or aptitude measure
\cite[pp.~233--234]{Khemlani2021Syllogistic}.

\begin{hbconvention}{def:p19-r33-lookup}{Selecting the relevant comparison}
For a validity question, translate the premises, preserve witness assumptions
and test all permitted situations; use \cref{thm:p19-r33-finite-check}
or the P2 calculus. For a predicted-answer question, state which model
version, parameter values, construction, response format and acceptance
mechanism are used. For a claim about people, retain the recruited sample,
repeated sessions, exclusions, item set and uncertainty.
An easy valid \(Iab,Abc\) item does not establish competence on
\(Eab,Abc\); one found counterexample does not establish NVC;
one high correlation does not identify a unique mechanism; and a belief-
content fit does not show transfer to neutral contents, clinical populations
or unrestricted natural-language arguments.
\end{hbconvention}
\index{syllogistic reasoning!psychological theories}
\index{mental models!mReasoner}\index{probabilistic representation model}
\index{figural effect}\index{syllogistic reasoning!existential import}
